% !TEX program = xelatex
% ============================================================
%  Arabic Course Syllabus Template — Nibras Center
%  Compile with: xelatex main.tex (twice)
% ============================================================
%  Features:
%    - A4, 11pt, article class
%    - Course info block (code, credits, prerequisites)
%    - Sections: description, objectives, outcomes,
%      weekly schedule, grading, policies, references
% ============================================================

\documentclass[11pt,a4paper]{article}

% ---- Geometry ----
\usepackage[a4paper,margin=2.5cm,top=2.5cm,bottom=2.5cm,headheight=15pt]{geometry}

% ---- Core packages (order matters for bidi) ----
\usepackage{url}
\usepackage{amsmath}
\usepackage{graphicx}
\usepackage{array}
\usepackage{booktabs}
\usepackage{tabularx}
\usepackage{multirow}
\usepackage{enumitem}
\usepackage{float}
\usepackage{caption}
\usepackage{titlesec}
\usepackage{fancyhdr}
\usepackage{setspace}
\usepackage[table]{xcolor}

% ---- RTL / Arabic language support ----
\usepackage{bidi}
\usepackage[no-math]{fontspec}
\usepackage[quiet]{polyglossia}

% ---- Fonts ----
\setmainfont[Scale=1]{NotoKufiArabic-Regular.ttf}
\newfontfamily\arabicfont[Script=Arabic,Scale=0.95]{NotoKufiArabic-Regular.ttf}
\newfontfamily\englishfont[Scale=0.95]{TeX Gyre Termes}
\setmonofont{NotoKufiArabic-Regular.ttf}

% ---- Languages ----
\setdefaultlanguage[calendar=gregorian,hijricorrection=1]{arabic}
\setotherlanguage[variant=british]{english}

% ---- Hyperref ----
\usepackage[breaklinks,hidelinks]{hyperref}
\hypersetup{
  pdftitle={خطة دورة},
  pdfauthor={اسم المدرّس}
}

% ---- Captions in Arabic ----
\addto\captionsarabic{%
  \renewcommand{\contentsname}{الفهرس}%
  \renewcommand{\figurename}{الشكل}%
  \renewcommand{\tablename}{الجدول}%
  \renewcommand{\refname}{المراجع}%
}

% ---- Section formatting ----
\titleformat{\section}{\Large\bfseries}{\thesection}{0.7em}{}
\titleformat{\subsection}{\large\bfseries}{\thesubsection}{0.6em}{}
\titlespacing*{\section}{0pt}{1.2ex plus 0.3ex}{0.8ex plus 0.2ex}
\titlespacing*{\subsection}{0pt}{0.8ex plus 0.2ex}{0.5ex plus 0.1ex}

% ---- Headers / footers ----
\pagestyle{fancy}
\fancyhf{}
\fancyhead[R]{\small\itshape خطة دورة: اسم المقرر}
\fancyhead[L]{\small اسم المدرّس}
\fancyfoot[C]{\small\thepage}
\renewcommand{\headrulewidth}{0.3pt}
\renewcommand{\footrulewidth}{0pt}

% ---- List spacing ----
\setlist[itemize]{topsep=0.3em, itemsep=0.15em, parsep=0pt}
\setlist[enumerate]{topsep=0.3em, itemsep=0.15em, parsep=0pt}

% ---- Table column types ----
\newcolumntype{L}[1]{>{\raggedright\arraybackslash}p{#1}}
\newcolumntype{C}[1]{>{\centering\arraybackslash}p{#1}}

% ---- Colors ----
\definecolor{headerColor}{HTML}{1A3C6B}
\definecolor{rowColor}{HTML}{F0F4F8}

% ---- Line spacing ----
\setstretch{1.4}

\begin{document}

% ============================================================
% TITLE BLOCK
% ============================================================
\begin{center}
{\LARGE\bfseries خطة دورة: اسم المقرر}\\[0.5em]
{\Large القسم، الجامعة}\\[1em]
\end{center}

% ---- Course info table ----
\begin{table}[H]
\centering
\renewcommand{\arraystretch}{1.4}
\begin{tabularx}{\textwidth}{L{4cm} X L{4cm} X}
\toprule
\textbf{رمز المقرر:} \LR{XXX 101} & \textbf{عدد الساعات:} 3 & \textbf{المستوى:} الأول & \textbf{الفصل:} --- \\
\textbf{المتطلبات السابقة:} --- & \textbf{المتطلبات المرافقة:} --- & \textbf{لغة التدريس:} العربية & \textbf{نوع المقرر:} إجباري \\
\bottomrule
\end{tabularx}
\end{table}

\vspace{1em}

% ---- Instructor info ----
\begin{table}[H]
\centering
\renewcommand{\arraystretch}{1.3}
\begin{tabularx}{\textwidth}{L{4cm} X L{4cm} X}
\toprule
\textbf{اسم المدرّس:} --- & \textbf{البريد:} \LR{---@---.edu} & \textbf{المكتب:} --- & \textbf{ساعات المكتب:} --- \\
\bottomrule
\end{tabularx}
\end{table}

\vspace{1em}

% ============================================================
% 1. COURSE DESCRIPTION
% ============================================================
\section{وصف المقرر}
\label{sec:description}

% TODO: اكتب وصفاً مختصراً للمقرر (3-5 جمل).
هذا نص تجريبي لوصف المقرر. يقدم نظرة عامة على محتوى المقرر ومجاله وأهميته في البرنامج الدراسي.

% ============================================================
% 2. OBJECTIVES
% ============================================================
\section{أهداف المقرر}
\label{sec:objectives}

% TODO: اكتب أهداف المقرر.
\begin{enumerate}
    \item فهم المفاهيم الأساسية في ...
    \item تطبيق المهارات العملية في ...
    \item تحليل المسائل المتعلقة بـ ...
    \item تقييم الحلول المختلفة لـ ...
\end{enumerate}

% ============================================================
% 3. LEARNING OUTCOMES
% ============================================================
\section{مخرجات التعلّم}
\label{sec:outcomes}

% TODO: اكتب مخرجات التعلّم المتوقعة.
بنهاية هذا المقرر، سيكون الطالب قادراً على:
\begin{enumerate}
    \item \textbf{المعرفة:} وصف ...
    \item \textbf{المهارة:} تطبيق ...
    \item \textbf{الاتجاه:} تقدير ...
\end{enumerate}

% ============================================================
% 4. WEEKLY SCHEDULE
% ============================================================
\section{الخطة الأسبوعية}
\label{sec:schedule}

% TODO: عدّل الخطة الأسبوعية حسب مقررك.
\begin{table}[H]
\centering
\small
\caption{الخطة الأسبوعية للمقرر}
\renewcommand{\arraystretch}{1.3}
\begin{tabularx}{\textwidth}{C{1.5cm} X C{2.5cm} C{2.5cm}}
\toprule
\textbf{الأسبوع} & \textbf{الموضوع} & \textbf{القراءة} & \textbf{التقييم} \\
\midrule
1 & مقدمة وتعريف بالمقرر & الفصل 1 & --- \\
2 & الموضوع الأول & الفصل 2 & واجب 1 \\
3 & الموضوع الثاني & الفصل 3 & --- \\
4 & الموضوع الثالث & الفصل 4 & اختبار 1 \\
5 & الموضوع الرابع & الفصل 5 & واجب 2 \\
6 & المراجعة والاختبار الشهري & --- & اختبار شهري \\
7 & الموضوع الخامس & الفصل 6 & --- \\
8 & الموضوع السادس & الفصل 7 & واجب 3 \\
9 & الموضوع السابع & الفصل 8 & اختبار 2 \\
10 & الموضوع الثامن & الفصل 9 & واجب 4 \\
11 & الموضوع التاسع & الفصل 10 & --- \\
12 & المراجعة العامة & --- & --- \\
13 & العروض التقديمية & --- & عرض \\
14 & الاختبار النهائي & --- & اختبار نهائي \\
\bottomrule
\end{tabularx}
\end{table}

% ============================================================
% 5. GRADING
% ============================================================
\section{نظام التقييم}
\label{sec:grading}

% TODO: عدّل نظام التقييم حسب مقررك.
\begin{table}[H]
\centering
\caption{توزيع الدرجات}
\begin{tabular}{|l|c|}
\hline
\textbf{البند} & \textbf{الدرجة} \\
\hline
الواجبات & 20\% \\
الاختبارات القصيرة & 15\% \\
الاختبار الشهري & 20\% \\
العرض التقديمي & 10\% \\
الاختبار النهائي & 35\% \\
\hline
\textbf{المجموع} & \textbf{100\%} \\
\hline
\end{tabular}
\end{table}

% ============================================================
% 6. POLICIES
% ============================================================
\section{السياسات الأكاديمية}
\label{sec:policies}

\subsection{الحضور والغياب}
% TODO: اكتب سياسة الحضور.
الحضور إجباري. لا يسمح بغياب أكثر من \LR{25\%} من المحاضرات، وإلا يُحرم الطالب من الاختبار النهائي.

\subsection{التسليم المتأخر}
% TODO: اكتب سياسة التسليم المتأخر.
يُخصم \LR{10\%} من درجة الواجب عن كل يوم تأخير.

\subsection{النزاهة الأكاديمية}
% TODO: اكتب سياسة النزاهة الأكاديمية.
يُتوقع من جميع الطلاب الالتزام بأعلى معايير النزاهة الأكاديمية. الانتحال أو الغش يؤدي إلى صفر في التقييم.

% ============================================================
% 7. REFERENCES
% ============================================================
\section{المراجع}
\label{sec:references}

\subsection{المراجع الرئيسية}
\begin{enumerate}
    \item Author, A. (2024). \emph{اسم الكتاب}. دار النشر.
\end{enumerate}

\subsection{المراجع المساندة}
\begin{enumerate}
    \item Author, B. (2023). \emph{اسم الكتاب}. دار النشر.
\end{enumerate}

\end{document}
